\section*{NS92: a direct arithmetic budget for the selected cost}

\paragraph{Scope and prior work.}
The continuous average and complete norm of its generator are NS61; the
selected Mobius/logarithmic direction is NS73, and the finite lower
numerator is NS91. We apply that average to signed coefficient tails,
obtaining an unconditional, computable upper cost and an endpoint-only
feasible update. No uniform bound on the arithmetic sum or cofinal gain
is proved. This is not a new RH criterion or new best finite error.

Use $\mathcal H=L^2((0,\infty),dt/t^2)$, $a_x(t)=A(t/x)$,
$A(z)=\int_0^z\{u\}\,du/u$, $\tau=\log_+t$, $R_N=\tau-\Pi_N\tau$,
and $E_N=\|R_N\|^2$. Put
\[
 d_n=-\mu(n)\log(2N/n)\quad(N<n\le2N),\quad
 \ell_N=\langle R_N,\sum d_na_n\rangle,\quad
 K_N=\|(I-\Pi_N)\sum d_na_n\|^2.
\]
The old coefficients of NS90 have no effect on $\ell_N$ or $K_N$.
Write $\alpha=\log2$ and $\kappa=\log(2\pi)-\gamma$.

\paragraph{Exact complete-space average.}
For $h(z)=A(z)-\{z\}$, NS61 proves the strong integral identity
\[
 \frac{d}{dx}(xa_x(t))=h(t/x),\quad
 h(t/x)=0\ (t<x),\quad \|h(\,\cdot/x)\|^2=\kappa/x.
\]
Define
\[
 T_N(x)=\sum_{x<n\le2N}d_n/n\quad(N\le x\le2N),\qquad
 F_N=\sum_{N<n\le2N}d_na_n-NT_N(N)a_N.
\]
Finite interchange gives the exact identity
\begin{equation}\label{ns92-average}
 F_N(t)=\int_N^{2N}T_N(x)h(t/x)\,dx.
\end{equation}
Each summand contributes $(d_n/n)\int_N^n h(t/x)dx
=d_na_n-(Nd_n/n)a_N$. Thus $F_N=0$ on $0<t<N$, its weighted
coefficient sum is zero, $\langle R_N,F_N\rangle=\ell_N$, and its
old-space projection complement is the original selected direction.

\paragraph{Proved complete cost inequalities.}
Minkowski and weighted Cauchy--Schwarz give
\begin{align}
 K_N\le\|F_N\|^2
 &\le U_N:=\kappa\left(\int_N^{2N}|T_N(x)|\frac{dx}{\sqrt x}\right)^2
 \label{ns92-cost}\\
 &\le\kappa\alpha S_N,\qquad S_N:=\int_N^{2N}T_N(x)^2dx.
 \label{ns92-square}
\end{align}
All norms are over the whole half-line. No physical cutoff, model cost,
truncated Gram, or omitted infinite tail enters these inequalities.
For integer $N$, they are finite sums:
\[
 S_N=\sum_{m=N}^{2N-1}T_N(m)^2,\qquad
 U_N=4\kappa\left(\sum_{m=N}^{2N-1}|T_N(m)|
                    (\sqrt{m+1}-\sqrt m)\right)^2.
\]
Form each signed Mobius suffix before taking its absolute value.
This is an upper bound, not an equality for the projected cost.

\paragraph{Constructive step without a new-block inverse.}
NS91 supplies a complete lower numerator $0<b_N\le\ell_N$ at the
tested sizes. For any specified $\lambda\ge0$,
\[
 E_N-\|R_N-\lambda F_N\|^2
 =2\lambda\ell_N-\lambda^2\|F_N\|^2
 \ge2\lambda b_N-\lambda^2U_N.
\]
The exact step $\lambda=b_N/U_N$ guarantees
$E_N-E_{2N}\ge b_N^2/U_N$. For a literal rational step the same
inequality applies directly; the certificates use $\lambda=1/512$.
Add $-\lambda NT_N(N)$ to the old coefficient at $N$, and
$\lambda d_n$ at each new index. No refit or new-block Gram inverse is
needed to execute the update. The starting optimized residual still uses
inherited old-space coefficient enclosures; no fast cofinal rule for them
is claimed.

\paragraph{What an eventual estimate would require.}
Let $M(x)=\sum_{n\le x}\mu(n)$ and
$B_N=\sup_{N\le x\le y\le2N}|M(y)-M(x)|$.
Partial summation with $f(x)=\log(2N/x)/x$, $f(2N)=0$, gives
\[
 T_N(x)=\int_x^{2N}(M(t)-M(x))f'(t)dt,\qquad
 |T_N(x)|\le B_N f(x).
\]
Consequently
\begin{equation}\label{ns92-oscillation}
 S_N\le (1-\alpha)^2B_N^2/N,\qquad
 U_N\le\kappa\alpha(1-\alpha)^2B_N^2/N,
\end{equation}
using $\int_1^2\log^2(2/u)\,du/u^2=(1-\alpha)^2$.
The elementary $B_N\le N$ yields only $O(N)$. A square-root block
oscillation estimate would give a bounded budget, but is neither proved
nor assumed. The hypothesis $|M(x)|\le Cx^\theta$ gives only
$U_N=O(N^{2\theta-1})$. An RH-dependent estimate, or an arbitrarily
small positive power, must not be replaced by a bounded cost.

If $b_N\ge\varepsilon E_N$ and $S_N\le C$ eventually on $N=2^j$,
with fixed $\varepsilon,C>0$, then NS83's eventual $E_N\ge c/\log N$
implies
\[
 \frac{E_N-E_{2N}}{E_N}\ge
 \frac{\varepsilon^2c}{\kappa\alpha^2 C}\frac1j.
\]
Both eventual hypotheses remain open. This is a sufficient implication,
not a necessity claim or a proof of the missing arithmetic theorem.
The projected cost may be far smaller than $U_N$; failure of this budget
would not exclude the original PR42 route.

\paragraph{Finite checks and their limits.}
At $N=16,64,256,1024,4096,16384,65536$, two-precision scalar sums give
$0.00098<U_N<0.00220$ on this finite list. At $N=256$, $U_N<0.001063$,
about $60.74$ times the inherited true projected cost. Combining it with
NS91 yields $b_N^2/(U_NE_N)>0.000464$. The cheap budget certifies relative gain $>0.000443$ for $\lambda=1/512$.
The separate complete physical calculation certifies gain $>0.0157$ for
$\lambda=1/16$; this larger step uses its sharper physical cost, not the
coarse scalar upper budget. Neither step refits the other old coefficients.
No finite percentage or upper bound is extrapolated.

Independent physical integration evaluates the old error, direction norm,
and pairing on every unit cell through $T$, then adds their complete main
tails and all main--remainder and remainder--remainder terms. For
$r_i=u_i\log t+v_i+\epsilon_i$, $|\epsilon_i|\le B_i/t$, set
$q_i=u_i\log T+v_i$, $M_i=(q_i^2+2u_iq_i+2u_i^2)/T$, and
$V_i=B_i^2/(3T^3)$. The tail pairing main term is
\[
 (q_1q_2+u_1q_2+u_2q_1+2u_1u_2)/T
\]
with radius $\sqrt{M_1V_2}+\sqrt{M_2V_1}+\sqrt{V_1V_2}$.
For a norm the radius is $2\sqrt{M_iV_i}+V_i$. The old residual's entire
exterior interval $0<t<1$ is retained. Replays use 256/384 bits and a
doubled physical cutoff with the inherited old-coefficient cutoff replay.
They establish a finite feasible decrease, not a new best error, eventual
gain, or proof of RH. Review is by the author only.
